%*****************************************
\chapter{Ergebnisse}\label{ch:results}
%*****************************************
\begin{temp}
Auch die Ergebnisse können in mehrere Kapitel aufgeteilt werden.
Falls die logische Struktur der Arbeit es anbietet, können auch mehrere Abfolgen von Methodik und Ergebnissen sinnvoll sein.
Beispiel: Ziel 1 ist die Erfassung von Nutzeranforderungen und Ziel 2 die Entwicklung oder Konzeption einer Software auf Basis der Anforderungen.
Wenn nun das Methodikkapitel für Ziel 1 eine Methodik für Interviews beschreibt und für Ziel 2 eine Softwareentwicklungsmethodik, die zu den Ergebnissen der Interviews passt, dann gehören die Interviewergebnisse ja inhaltlich dazwischen.
%Wichtig ist, dass man zwischen Ausführung und Ergebnis unterscheiden kann.

\section{Zielerreichung}
In der Zielerreichung wird vor allem zusammenfassend dargestellt, welche Ergebnisse zu den formulierten Zielen erreicht wurden.
Bei der Formulierung sollten wichtige Schlüsselbegriffe verwendet werden sowie explizit und strukturiert jedes Ziel aus der Einleitung aufgegriffen werden.
%Wurden Fragen formuliert, können auch diese hier beantwortet werden. % Bei uns gibt es keine Forschungsfragen sondern Aufgaben und Ziele
So soll es möglich sein, dass ein Leser von der Arbeit lediglich die Einleitung und die Zielerreichung liest und doch die Ergebnisse der Arbeit erfassen kann.
\end{temp}
